% =============================================================================
% Module 1e: Solutions to Exercises
% =============================================================================
% This file is conditionally included based on \ifsolutions
% =============================================================================

\section{Solutions to Exercises}
\label{sec:linearized_solutions}

\noindent
Full worked solutions are also available in the Mathematica script
\solnlink{01e_Linearized_Gravity/problems\_01e.wl}{problems\_01e.wl},
which verifies each answer symbolically.

% ----- Exercise 1 -----
\subsection*{Exercise~\ref{ex:newtonian_limit_prob}: Newtonian Limit from the Geodesic Equation}

\textbf{(a)} The Christoffel symbols to first order in $h$ are
(eq.~\eqref{eq:christoffel_linearized}):
\[
    \christoffel{\alpha}{\mu}{\nu}
    = \half \eta^{\alpha\beta}\bigl(
    \partial_\mu h_{\nu\beta} + \partial_\nu h_{\mu\beta}
    - \partial_\beta h_{\mu\nu}\bigr) \,.
\]
For $\mu = \nu = 0$ and $\alpha = i$ (spatial):
\[
    \christoffel{i}{0}{0}
    = \half \eta^{i\beta}\bigl(
    \partial_0 h_{0\beta} + \partial_0 h_{0\beta}
    - \partial_\beta h_{00}\bigr) \,.
\]
Using the static-field condition $\partial_0 h_{\mu\nu} = 0$,
the first two terms vanish.  Since $\eta^{ij} = \delta^{ij}$ and
$\eta^{i0} = 0$:
\[
    \boxed{\christoffel{i}{0}{0} = -\half \partial_i h_{00}}
\]

\textbf{(b)} The geodesic equation for spatial components is:
\[
    \frac{d^2 x^i}{d\tau^2}
    + \christoffel{i}{\alpha}{\beta}\,
    \frac{dx^\alpha}{d\tau}\, \frac{dx^\beta}{d\tau} = 0 \,.
\]
In the slow-motion limit, $|dx^j/d\tau| \ll c\,dt/d\tau$, so we keep
only the $\alpha = \beta = 0$ term:
\[
    \frac{d^2 x^i}{d\tau^2}
    + \christoffel{i}{0}{0}\left(\frac{dx^0}{d\tau}\right)^{\!2} = 0 \,.
\]
Substituting $\christoffel{i}{0}{0} = -\half \partial_i h_{00}$ and
$dx^0/d\tau = c\, dt/d\tau$:
\[
    \frac{d^2 x^i}{d\tau^2}
    = \half c^2 \partial_i h_{00}\left(\frac{dt}{d\tau}\right)^{\!2} .
\]
Since $d\tau \approx dt$ for slow motion ($\gamma \approx 1$):
\[
    \boxed{\frac{d^2 x^i}{dt^2} = \half c^2\, \partial_i h_{00}}
\]

\textbf{(c)} Comparing with Newton's law $d^2 x^i/dt^2 = -\partial_i \Phi$:
\[
    -\partial_i \Phi = \half c^2\, \partial_i h_{00}
    \quad\Longrightarrow\quad
    \boxed{h_{00} = -\frac{2\Phi}{c^2}}
\]

% ----- Exercise 2 -----
\subsection*{Exercise~\ref{ex:weak_field_schwarzschild}: Weak-Field Schwarzschild}

\textbf{(a)} The Schwarzschild metric has
$g_{rr} = (1 - R_S/r)^{-1}$.  Expanding to first order in
$\epsilon \equiv R_S/r \ll 1$:
\[
    (1 - \epsilon)^{-1} = 1 + \epsilon + \epsilon^2 + \cdots
    \approx 1 + \frac{R_S}{r} \,.
\]
So the metric becomes:
\[
    ds^2 \approx -\left(1 - \frac{R_S}{r}\right) c^2 dt^2
    + \left(1 + \frac{R_S}{r}\right) dr^2
    + r^2\, d\Omega^2 \,.
\]

\textbf{(b)} In Cartesian-like coordinates ($dx^2 + dy^2 + dz^2 =
dr^2 + r^2 d\Omega^2$ to zeroth order), and noting that the
distinction between $dr^2 + r^2 d\Omega^2$ and $(1 + R_S/r)(dr^2 +
r^2 d\Omega^2)$ versus $(1 + R_S/r) dr^2 + r^2 d\Omega^2$ only
matters at second order in $R_S/r$:
\[
    ds^2 \approx -\left(1 - \frac{R_S}{r}\right) c^2 dt^2
    + \left(1 + \frac{R_S}{r}\right)(dx^2 + dy^2 + dz^2) \,.
\]
With $\Phi = -GM/r$ and $R_S = 2GM/c^2$:
\[
    1 - \frac{R_S}{r} = 1 + \frac{2\Phi}{c^2} \,, \qquad
    1 + \frac{R_S}{r} = 1 - \frac{2\Phi}{c^2} \,.
\]
This is exactly the weak-field metric~\eqref{eq:weak_field_metric}.

\textbf{(c)} We identify $h_{00} = -2\Phi/c^2 = R_S/r$ and
$h_{ij} = -2\Phi/c^2\, \delta_{ij} = (R_S/r)\, \delta_{ij}$.
Indeed $h_{00} = h_{ij}/\delta_{ij}$, confirming that the time and
space perturbations are equal in magnitude.  This equality is the
origin of the factor-of-two in light bending.

% ----- Exercise 3 -----
\subsection*{Exercise~\ref{ex:refractive_index_prob}: Effective Refractive Index}

The effective refractive index is $n - 1 = 2GM/(c^2 r)$.
Using $G = 6.674 \times 10^{-11}$~m$^3$\,kg$^{-1}$\,s$^{-2}$,
$M_\odot = 1.989 \times 10^{30}$~kg, $c = 2.998 \times 10^8$~m/s:

\textbf{(a)} $r = R_\odot = 6.96 \times 10^8$~m:
\[
    n - 1 = \frac{2 \times 6.674\!\times\!10^{-11} \times 1.989\!\times\!10^{30}}
    {(2.998\!\times\!10^{8})^2 \times 6.96\!\times\!10^{8}}
    = \boxed{4.24 \times 10^{-6}}
\]
The weak-field approximation is excellent ($n - 1 \sim 10^{-6}$).

\textbf{(b)} $r = 1$~AU $= 1.496 \times 10^{11}$~m:
\[
    n - 1 = \frac{2GM_\odot}{c^2 \times 1.496\!\times\!10^{11}}
    = \boxed{1.97 \times 10^{-8}}
\]
Even weaker field; the Sun's effect at Earth's distance is tiny.

\textbf{(c)} $r = 1$~kpc $= 3.086 \times 10^{19}$~m:
\[
    n - 1 = \frac{2GM_\odot}{c^2 \times 3.086\!\times\!10^{19}}
    = \boxed{9.58 \times 10^{-17}}
\]
Completely negligible.  At kiloparsec scales, the gravitational
potential of a single solar mass is immeasurably weak; only the
cumulative effect of $\sim 10^{11}$--$10^{12}$ stars in a galaxy
produces observable lensing.

% ----- Exercise 4 -----
\subsection*{Exercise~\ref{ex:deflection_angle_derivation}: Deflection Angle from Null Geodesics}

The unperturbed photon path is $\mathbf{x}(l) = (l, b, 0)$, where
$l$ is the coordinate along the path and $b$ is the impact parameter.

\textbf{(a)} The Newtonian potential is $\Phi = -GM/r = -GM/\sqrt{l^2 + b^2}$.
The perpendicular gradient (in the $y$-direction) is:
\[
    \nabla_\perp \Phi = \frac{\partial \Phi}{\partial b}
    = \frac{GM\, b}{(l^2 + b^2)^{3/2}} \,.
\]

\textbf{(b)} The integral is:
\[
    \int_{-\infty}^{\infty} \frac{b\, dl}{(l^2 + b^2)^{3/2}}
    = \frac{b}{b^2} \left[\frac{l}{\sqrt{l^2 + b^2}}\right]_{-\infty}^{+\infty}
    = \frac{1}{b}(1 - (-1)) = \frac{2}{b} \,.
\]

\textbf{(c)} The deflection integral
(eq.~\eqref{eq:deflection_integral}) uses $\nabla_\perp \Phi$ to denote
the gradient of $|\Phi|$ perpendicular to the path.  The factor of~4
already accounts for both the time ($h_{00}$) and space ($h_{ij}$)
contributions.  The key subtlety is that the standard 3D line-of-sight
integral relates to the 2D projected formula as follows.

Using the projected (2D) lens-plane formulation
(Schneider, Ehlers \& Falco 1992; Congdon \& Keeton eq.~3.78), the
deflection by a surface mass density $\Sigma$ is:
\[
    \boldsymbol{\hat{\alpha}} = \frac{4G}{c^2}
    \int \frac{(\vecxi - \vecxi')}{|\vecxi - \vecxi'|^2}\,
    \Sigma(\vecxi')\, d^2\xi' \,.
\]
For a point mass projected at the origin,
$\Sigma = M\,\delta^{(2)}(\vecxi')$, and the integral gives directly:
\[
    \boxed{\alphahat = \frac{4GM}{c^2 b}}
\]
This is verified symbolically in
\mathlink{01e_Linearized_Gravity/weak_field_metric.wl}{weak\_field\_metric.wl}.

Equivalently, using the PPN formalism (Weinberg 1972; Will 1993),
the deflection is:
\[
    \alphahat = \frac{1 + \gamma_{\mathrm{PPN}}}{c^2}
    \int_{-\infty}^{\infty} \nabla_\perp |\Phi|\, dl \,,
\]
where $\gamma_{\mathrm{PPN}} = 1$ in GR (the spatial curvature
contribution).  With $\int \nabla_\perp |\Phi|\, dl = 2GM/b$ and
$(1 + 1) = 2$, this gives $\alphahat = (2/c^2)(2GM/b) = 4GM/(c^2 b)$,
confirming the result.

\textbf{(d)} For the Sun at the solar limb:
\begin{align*}
    \alphahat &= \frac{4 \times 6.674\!\times\!10^{-11} \times
    1.989\!\times\!10^{30}}{(2.998\!\times\!10^{8})^2 \times
    6.96\!\times\!10^{8}} \\
    &= 8.49 \times 10^{-6}~\text{rad}
    = 8.49 \times 10^{-6} \times \frac{180 \times 3600}{\pi}
    = \boxed{1.75''}
\end{align*}
This is the value confirmed by Eddington's 1919 expedition.

\textbf{(Galaxy lens.)} For $M = 10^{12}\,\Msun$ at $b = 10$~kpc
$= 3.086 \times 10^{20}$~m:
\[
    \alphahat = \frac{4 \times 6.674\!\times\!10^{-11} \times
    10^{12} \times 1.989\!\times\!10^{30}}
    {(2.998\!\times\!10^{8})^2 \times 3.086\!\times\!10^{20}}
    = 1.91 \times 10^{-5}~\text{rad}
    = \boxed{3.94''}
\]

% ----- Exercise 5 -----
\subsection*{Exercise~\ref{ex:superposition_prob}: Superposition of Deflections}

\textbf{(a)} Mass $M_1$ at $(-d/2, 0)$ and mass $M_2$ at $(+d/2, 0)$.
The photon passes at $\vecxi = (0, b)$.  The displacement vectors are:
\begin{align*}
    \vecxi - \vecxi_1 &= (d/2,\, b) \,, &
    |\vecxi - \vecxi_1|^2 &= d^2/4 + b^2 \,, \\
    \vecxi - \vecxi_2 &= (-d/2,\, b) \,, &
    |\vecxi - \vecxi_2|^2 &= d^2/4 + b^2 \,.
\end{align*}
The individual deflection angle vectors are:
\begin{align*}
    \boldsymbol{\hat{\alpha}}_1 &= \frac{4GM}{c^2}
    \frac{(d/2, b)}{d^2/4 + b^2} \,, \\
    \boldsymbol{\hat{\alpha}}_2 &= \frac{4GM}{c^2}
    \frac{(-d/2, b)}{d^2/4 + b^2} \,.
\end{align*}

\textbf{(b)} Adding the vectors:
\[
    \hat{\alpha}_x = \frac{4GM}{c^2(d^2/4 + b^2)}
    \left(\frac{d}{2} - \frac{d}{2}\right) = 0
\]
(the $x$-components cancel by symmetry).  The $y$-component:
\[
    \boxed{\hat{\alpha}_y = \frac{4GM}{c^2(d^2/4 + b^2)} \cdot 2b
    = \frac{8GMb}{c^2(d^2/4 + b^2)}}
\]

\textbf{(c)} In the limit $d \to 0$:
\[
    \hat{\alpha}_y \to \frac{8GMb}{c^2 b^2}
    = \frac{8GM}{c^2 b}
    = \frac{4G(2M)}{c^2 b}
\]
This is the deflection angle for a single mass of $2M$, as expected
from superposition when the two masses coincide.
